\section{Discussion}
\label{sec:discussion}

\subsection{Legal Theory: CSS as a Litigation Tool}
\label{sec:legal}

The Census Stability Score has a precise legal application in
post-\emph{Rucho} partisan gerrymandering litigation.

\paragraph{The constitutional argument.}
Under \emph{Rucho v.\ Common Cause} (2019), federal courts cannot adjudicate
partisan gerrymandering claims.
But state courts, applying their own constitutions, retain this jurisdiction.
The key factual question in state-court litigation is: \emph{does the enacted
map deviate from the natural map, and if so, can the deviation be explained
by legitimate geographic factors?}

A GeoSection map with CSS $\geq 0.90$ provides the strongest possible answer
to this question.
The map is the natural map.
It has been the natural map across three independent decennial censuses.
Any deviation from it---any enacted map that differs from the GeoSection
optimal---is \emph{not} the natural map and requires an explanation beyond
geography.

\begin{proposition}[CSS as Geographic Burden Shift]
\label{prop:burden-shift}
If a state has a GeoSection map with CSS$_{3\text{yr}} \geq 0.90$, then
for any enacted congressional map that differs from the GeoSection map,
the Legislature cannot justify the deviation on geographic grounds alone,
because the geographic structure of the state has consistently implied a
different partition across three independent census cycles.
\end{proposition}

The strength of Proposition~\ref{prop:burden-shift} increases with:
(a) the number of census years confirming the same outcome,
(b) the magnitude of the proportionality gap in the enacted map relative
to the GeoSection map, and
(c) the stability of the surrounding states (if all neighbouring states
have CSS $\geq 0.90$, the geographic argument is regional, not idiosyncratic).

\paragraph{The ``before living memory'' standard.}
A map with CSS $\geq 0.90$ across 2000, 2010, and 2020 \emph{existed before
living memory of the current partisan alignment}.
In 2000, the specific partisan geography that characterises today's Republican
dominance in exurban and rural areas was less pronounced; the urban-rural
divide was real but less extreme.
The fact that GeoSection produces the same 5\,D/9\,R outcome in 2000 as in
2020 for a state like Wisconsin means the outcome is not contingent on the
current partisan distribution.
It predates the 2016 political realignment.
It predates the post-2010 Tea Party Republican gains in suburban areas.
It is geography.

\paragraph{The Lorenz-stable map as a statute design tool.}
A proposed Districting Integrity Act could incorporate CSS as a procedural
threshold:
\begin{enumerate}
\item Any proposed congressional map must be accompanied by its CSS score
  computed against the GeoSection-optimal map.
\item Maps with CSS $\geq 0.90$ may be adopted by administrative certification
  (no legislative override required).
\item Maps with CSS $0.70$--$0.90$ require a majority-vote legislative approval
  with public disclosure of the CSS components.
\item Maps with CSS $< 0.70$ trigger a heightened judicial review process,
  because the proposed map deviates substantially from the state's durable
  geographic structure.
\end{enumerate}

This framework respects democratic self-governance: legislatures retain the
ability to adopt non-GeoSection maps, but they must publicly acknowledge the
degree to which their map departs from the natural geographic partition.

\paragraph{States with high CSS are the strongest litigation targets.}
When a high-CSS state has an enacted map that differs from GeoSection,
the discrepancy is maximally hard to explain by geography.
Iowa (CSS $\approx 0.95$, predicted) has a 2:2 equal first-level split that
reflects its near-uniform agricultural-density geography.
If Iowa's enacted map had a 1:3 urban peel of Des Moines, the discrepancy
could not be attributed to any geographic force: the algorithm, across three
censuses, has never selected a 1:3 peel for Iowa.

\subsection{Population Dynamics and Geographic Reorganisation}
\label{sec:pop-dynamics}

States with low CSS are not failed cases for GeoSection.
They are states where the population-geography relationship has genuinely
reorganised between census cycles.

\paragraph{The Sun Belt restructuring hypothesis.}
Arizona's Phoenix metro grew from approximately 2.8M people in 2000 to
4.9M in 2020---a 75\,\% increase in twenty years.
This growth extended the Phoenix metro boundary into the far East Valley,
the West Valley, and North Scottsdale, creating new geographic regions that
are contiguous with Phoenix's compact urban core but have distinct
socioeconomic and demographic characteristics.
For GeoSection, this means the isoperimetric boundary of ``the Phoenix urban
core'' has expanded: in 2000, the natural partition may have peeled a compact
2.8M-person metro from a 5.1M total state population ($p^* \approx 0.13$, $1:7$);
in 2020, the metro is 4.9M of 7.2M total ($p^* \approx 0.11$, $1:8$).
The ratio shift from $1:7$ to $1:8$ is not an algorithm failure.
It is a signal that Arizona's human geography has reorganised.

\paragraph{Declining states: stability through contraction.}
Ironically, states with declining populations may exhibit high CSS because
geographic concentration \emph{decreases} over time.
West Virginia, Mississippi, and Alabama have seen substantial rural
depopulation, with population concentrating in smaller areas.
This tends to sharpen the geographic features that GeoSection detects,
potentially increasing the isoperimetric gap between ratios and making the
natural ratio more robust.
We predict these states will have high CSS despite (or because of)
demographic decline.

\paragraph{The Lorenz drift proxy as an early warning.}
States where $p^*$ shifts between census years are states where the
population-geography Lorenz curve has restructured.
Monitoring $\Delta p^*$ between the ACS 5-year estimates (available annually)
and the decennial census provides an early warning of potential ratio shifts.
A state where $\Delta p^*$ exceeds $0.03$ between two ACS estimates should
flag for updated GeoSection analysis, potentially before the decennial
redistricting trigger.

\subsection{Limitations}

\paragraph{Tract alignment across census years.}
The Jaccard similarity computation requires mapping district assignments from
2000 and 2010 tract boundaries to 2020 boundaries.
Our spatial intersection method (population-weighted allocation to 2020 tracts)
introduces approximation error in areas where the 2000 tract network differs
significantly from 2020.
Urban cores, where tracts are frequently split, have the highest
approximation error.
This means Jaccard similarity is slightly underestimated in urban areas,
potentially underestimating CSS for states with large cities.

\paragraph{Population-only perturbation requires common graph.}
The Type I/Type II decomposition requires running GeoSection on the 2020
graph with 2000 populations.
This is an approximation: 2000 tract population totals must be allocated to
2020 tracts by spatial intersection, which introduces the same approximation
error as the Jaccard computation.
For the stability \emph{difference} (full cross-year minus population-only
perturbation), the approximation errors partially cancel; but they do not
cancel exactly.

\paragraph{Seed count for large states.}
California ($k=52$), Texas ($k=38$), and Florida ($k=28$) have uncertain
natural ratios at 50 seeds per ratio.
Cross-census stability for these large states requires higher seed counts
(200+ per ratio) to ensure convergence.
We report their 2020 natural ratios as provisional and exclude them from the
CSS computation until higher-confidence estimates are available.

\paragraph{2010 sweep not yet complete.}
The three-census CSS requires 2010 data.
The 2010 sweep is in progress.
Two-census CSS covers the largest twenty-year span (2000--2020) and is
theoretically the stronger argument (a shorter interval would not capture
the 2010 redistricting cycle's partisan context).
Three-census CSS will strengthen the claim by adding an intermediate
data point that lies \emph{between} the bookend years.

\paragraph{Congressional apportionment changes the comparison baseline.}
Fourteen states changed seat count between 2000 and 2020.
Our normalised ratio comparison addresses this, but the comparison is
intrinsically less direct than for states with fixed seat counts.
A state like Texas ($k=30$ in 2000, $k=38$ in 2020) has a ratio comparison
that conflates algorithm stability with apportionment-change effects.
The CSS for seat-count-changing states should be interpreted with this caveat.

\subsection{The 2D Stability Picture: (seed, census)}

Combining B.7's seed stability results with B.15's census stability predictions,
we can sketch a $2\times 2$ stability matrix:

\begin{table}[h]
\centering
\caption{$2\times 2$ stability matrix: seed stability (B.7) $\times$ census
  stability (B.15 predicted). Representative states in each cell.}
\label{tab:2d-stability}
\begin{tabular}{l l l}
\toprule
                  & Low census stability & High census stability \\
\midrule
High seed stability & AZ, NV, CO         & IA, KS, WI, IL, PA \\
Low seed stability  & TX, FL, CA         & (none expected) \\
\bottomrule
\end{tabular}
\end{table}

States in the top-right cell (high seed stability, high census stability)
have the strongest geographic determinism claim and are the most legally
defensible.
States in the bottom-right cell are theoretically expected to be empty:
if the geographic structure is durable across censuses, it should also be
easily detectable by the algorithm (high seed convergence).
States in the bottom-left cell (TX, FL, CA) are algorithmically difficult:
their geographic structure is complex enough that multiple equally-good
solutions exist within a single census year, \emph{and} the optimal solution
shifts across census years as population redistributes.
These states require the most procedural justification for any particular map.

\subsection{Predictive Model for CSS}

The theoretical framework suggests a predictive regression:
\[
  \mathrm{CSS}(s) \sim f(\text{pop\_growth\_rate},\; \text{urban\_boundary\_growth},\;
                        \text{seat\_count\_change},\; p^*_{2020})
\]
with expected signs:
$\partial\mathrm{CSS}/\partial\text{pop\_growth} < 0$ (faster growth
$\to$ lower CSS),
$\partial\mathrm{CSS}/\partial\text{boundary\_growth} < 0$ (expanding
city $\to$ Lorenz drift),
$\partial\mathrm{CSS}/\partial\text{seat\_change} < 0$ (apportionment
change $\to$ ratio comparison noise),
$\partial\mathrm{CSS}/\partial p^*_{2020}$ depends on whether the ratio
is in the metropolitan-fixity region.

This regression, fit on the complete three-census CSS dataset once available,
would allow CSS prediction for future census cycles before running the full
50-state sweep.
